\documentclass{amsart}
% For dealing with references we use the comment environment
\usepackage{verbatim}
\newenvironment{reference}{\comment}{\endcomment}
%\newenvironment{reference}{}{}
\newenvironment{slogan}{\comment}{\endcomment}
\newenvironment{history}{\comment}{\endcomment}

% For commutative diagrams we use Xy-pic
\usepackage[all]{xy}

% We use 2cell for 2-commutative diagrams.
\xyoption{2cell}
\UseAllTwocells

% We use multicol for the list of chapters between chapters
\usepackage{multicol}

% This is generally recommended for better output
\usepackage{lmodern}
\usepackage[T1]{fontenc}

% For cross-file-references

% Package for hypertext links:
\usepackage{hyperref}

% For any local file, say "hello.tex" you want to link to please

% Theorem environments.
%
\theoremstyle{plain}
\newtheorem{theorem}[subsection]{Theorem}
\newtheorem{proposition}[subsection]{Proposition}
\newtheorem{lemma}[subsection]{Lemma}

\theoremstyle{definition}
\newtheorem{definition}[subsection]{Definition}
\newtheorem{example}[subsection]{Example}
\newtheorem{exercise}[subsection]{Exercise}
\newtheorem{situation}[subsection]{Situation}

\theoremstyle{remark}
\newtheorem{remark}[subsection]{Remark}
\newtheorem{remarks}[subsection]{Remarks}

\numberwithin{equation}{subsection}

% Macros
%
\def\lim{\mathop{\mathrm{lim}}\nolimits}
\def\colim{\mathop{\mathrm{colim}}\nolimits}
\def\Spec{\mathop{\mathrm{Spec}}}
\def\Hom{\mathop{\mathrm{Hom}}\nolimits}
\def\Ext{\mathop{\mathrm{Ext}}\nolimits}
\def\SheafHom{\mathop{\mathcal{H}\!\mathit{om}}\nolimits}
\def\SheafExt{\mathop{\mathcal{E}\!\mathit{xt}}\nolimits}
\def\Sch{\mathit{Sch}}
\def\Mor{\mathop{\mathrm{Mor}}\nolimits}
\def\Ob{\mathop{\mathrm{Ob}}\nolimits}
\def\Sh{\mathop{\mathit{Sh}}\nolimits}
\def\NL{\mathop{N\!L}\nolimits}
\def\CH{\mathop{\mathrm{CH}}\nolimits}
\def\proetale{{pro\text{-}\acute{e}tale}}
\def\etale{{\acute{e}tale}}
\def\QCoh{\mathit{QCoh}}
\def\Ker{\mathop{\mathrm{Ker}}}
\def\Im{\mathop{\mathrm{Im}}}
\def\Coker{\mathop{\mathrm{Coker}}}
\def\Coim{\mathop{\mathrm{Coim}}}

% Boxtimes
%
\DeclareMathSymbol{\boxtimes}{\mathbin}{AMSa}{"02}

%
% Macros for moduli stacks/spaces
%
\def\QCohstack{\mathcal{QC}\!\mathit{oh}}
\def\Cohstack{\mathcal{C}\!\mathit{oh}}
\def\Spacesstack{\mathcal{S}\!\mathit{paces}}
\def\Quotfunctor{\mathrm{Quot}}
\def\Hilbfunctor{\mathrm{Hilb}}
\def\Curvesstack{\mathcal{C}\!\mathit{urves}}
\def\Polarizedstack{\mathcal{P}\!\mathit{olarized}}
\def\Complexesstack{\mathcal{C}\!\mathit{omplexes}}
% \Pic is the operator that assigns to X its picard group, usage \Pic(X)
% \Picardstack_{X/B} denotes the Picard stack of X over B
% \Picardfunctor_{X/B} denotes the Picard functor of X over B
\def\Pic{\mathop{\mathrm{Pic}}\nolimits}
\def\Picardstack{\mathcal{P}\!\mathit{ic}}
\def\Picardfunctor{\mathrm{Pic}}
\def\Deformationcategory{\mathcal{D}\!\mathit{ef}}

\setlength{\emergencystretch}{3em}
\begin{document}
\title[Stacks Project, Tag 079P]{1344 --- Every unbounded complex in a Grothendieck abelian category admits a functorial termwise-injective K-injective resolution}
\maketitle
\noindent Copyright (C) 2005--2025 Johan de Jong. Stacks Project authors. Source: \href{https://stacks.math.columbia.edu/tag/079P}{Tag 079P}.

\noindent Editorial difficulty note (not author text). Difficulty H2: The construction alternates functorial homotopy-killing pushouts with an injective-term cone enlargement along an ordinal of uniformly large cofinality. A bounded-size acyclic test class, transfinite acyclic filtration and exact inverse Hom limits control all unbounded complexes; this is well beyond ordinary qualification resolution theory..

\noindent Editorial provenance note (not author text). History: excerpt from revision \texttt{a0444\allowbreak 6e57e\allowbreak c1fbc\allowbreak 252a8\allowbreak 71afc\allowbreak ec775\allowbreak 2fb28\allowbreak 07b14}, injectives.tex, lines 1806--1887. Reference presentation was adapted; original-author context and core support blocks were retained or added. The mathematical wording is unchanged. Licensed under GNU FDL 1.2 or later; no invariant sections or cover texts. See ../licenses/COPYING. Context blocks reproduce author definitions and supporting results; they are not additional counted problems.

\begin{definition}
\label{definition-grothendieck-conditions}
Let $\mathcal{A}$ be an abelian category. We name some conditions
\begin{enumerate}
\item[AB3] $\mathcal{A}$ has direct sums,
\item[AB4] $\mathcal{A}$ has AB3 and direct sums are exact,
\item[AB5] $\mathcal{A}$ has AB3 and filtered colimits are exact.
\end{enumerate}
Here are the dual notions
\begin{enumerate}
\item[AB3*] $\mathcal{A}$ has products,
\item[AB4*] $\mathcal{A}$ has AB3* and products are exact,
\item[AB5*] $\mathcal{A}$ has AB3* and cofiltered limits are exact.
\end{enumerate}
We say an object $U$ of $\mathcal{A}$ is a {\it generator} if
for every $N \subset M$, $N \not = M$ in $\mathcal{A}$ there exists a morphism
$U \to M$ which does not factor through $N$.
We say $\mathcal{A}$ is a {\it Grothendieck abelian category} if
it has AB5 and a generator.
\end{definition}

\begin{definition}
\label{definition-size}
Let $\mathcal{A}$ be a Grothendieck abelian category.
Let $M$ be an object of $\mathcal{A}$.
The {\it size} $|M|$ of $M$ is the cardinality of the set of subobjects
of $M$.
\end{definition}

\begin{definition}
\label{definition-small}
Let $\mathcal{C}$ be a category, let $I \subset \text{Arrows}(\mathcal{C})$,
and let $\alpha$ be an ordinal. An object $A$ of $\mathcal{C}$ is said to
be {\it $\alpha$-small with respect to $I$} if whenever $\{B_\beta\}$ is
a system over $\alpha$ with transition maps in $I$, then
the map (\ref{equation-compare}) is an isomorphism.
\end{definition}

\begin{definition}
\label{derived-definition-K-injective}
Let $\mathcal{A}$ be an abelian category. A complex $I^\bullet$
is {\it K-injective} if for every acyclic complex $M^\bullet$ we
have $\Hom_{K(\mathcal{A})}(M^\bullet, I^\bullet) = 0$.
\end{definition}

\medskip\noindent
We begin with a few set theoretic remarks.
Let $\{B_{\beta}\}_{\beta \in \alpha}$ be an inductive system of
objects in some category $\mathcal{C}$, indexed by
an ordinal $\alpha$. Assume that $\colim_{\beta \in \alpha} B_\beta$
exists in $\mathcal{C}$. If $A$ is an object of $\mathcal{C}$, then there is a
natural map
\begin{equation}
\label{equation-compare}
\colim_{\beta \in \alpha} \Mor_\mathcal{C}(A, B_\beta)
\longrightarrow
\Mor_\mathcal{C}(A, \colim_{\beta \in \alpha} B_\beta).
\end{equation}
because if one is given a map $A \to B_\beta$ for some $\beta$, one
naturally gets a map from $A$  into the colimit by composing with
$B_\beta \to \colim_{\beta \in \alpha} B_\alpha$.
Note that the left colimit is one of sets! In general, (\ref{equation-compare})
is neither injective or surjective.

\begin{lemma}
\label{lemma-surjection-bounded-size}
Let $\mathcal{A}$ be a Grothendieck abelian category with generator $U$.
Let $c$ be the function on cardinals defined by
$c(\kappa) = |\bigoplus_{\alpha \in \kappa} U|$. If $\pi : M \to N$ is a
surjection then there exists a subobject $M' \subset M$ which surjects
onto $N$ with $|M'| \leq c(|N|)$.
\end{lemma}

\begin{proof}
For every proper subobject $N' \subset N$ choose a morphism
$\varphi_{N'} : U \to M$ such that $U \to M \to N$ does not factor
through $N'$. Set
$$
M' = \Im\left(
\bigoplus\nolimits_{N' \subset N} \varphi_{N'} :
\bigoplus\nolimits_{N' \subset N} U \longrightarrow M\right)
$$
Then $M'$ works.
\end{proof}


\begin{lemma}
\label{lemma-acyclic-quotient-complexes-bounded-size}
Let $\mathcal{A}$ be a Grothendieck abelian category. There exists a cardinal
$\kappa$ such that given any acyclic complex $M^\bullet$ we have
\begin{enumerate}
\item if $M^\bullet$ is nonzero, there is a nonzero subcomplex
$N^\bullet$ which is bounded above, acyclic, and $|N^n| \leq \kappa$,
\item there exists a surjection of complexes
$$
\bigoplus\nolimits_{i \in I} M_i^\bullet \longrightarrow M^\bullet
$$
where $M_i^\bullet$ is bounded above, acyclic, and $|M_i^n| \leq \kappa$.
\end{enumerate}
\end{lemma}

\begin{proof}
Choose a generator $U$ of $\mathcal{A}$. Denote $c$ the function of
Lemma \ref{lemma-surjection-bounded-size}.
Set $\kappa = \sup \{c^n(|U|), n = 1, 2, 3, \ldots\}$.
Let $n \in \mathbf{Z}$ and let $\psi : U \to M^n$ be a morphism.
In order to prove (1) and (2) it suffices to prove there exists a subcomplex
$N^\bullet \subset M^\bullet$ which is bounded above, acyclic, and
$|N^m| \leq \kappa$, such that $\psi$ factors through $N^n$.
To do this set $N^n = \Im(\psi)$, $N^{n + 1} = \Im(U \to M^n \to M^{n + 1})$,
and $N^m = 0$ for $m \geq n + 2$.
Suppose for some $k \leq n$ we have constructed $N^m \subset M^m$
for all $n > m \geq k$ such that
\begin{enumerate}
\item $\text{d}(N^m) \subset N^{m + 1}$ for all $m \geq k$,
\item $\Im(N^{m - 1} \to N^m) = \Ker(N^m \to N^{m + 1})$ for
all $m \geq k + 1$, and
\item $|N^m| \leq c^{\max\{n - m, 0\}}(|U|)$ for all $m \geq k$.
\end{enumerate}
Because $M^\bullet$ is acyclic, we see that the subobject
$\text{d}^{-1}(\Ker(N^k \to N^{k + 1})) \subset M^{k - 1}$ surjects onto
$\Ker(N^k \to N^{k + 1})$. Thus we can choose $N^{k - 1} \subset M^{k - 1}$
surjecting onto $\Ker(N^k \to N^{k + 1})$ with
$|N^{k - 1}| \leq c^{n - k + 1}(|U|)$ by
Lemma \ref{lemma-surjection-bounded-size}. The proof is finished by
induction on $k$.
\end{proof}


\begin{lemma}
\label{lemma-characterize-K-injective}
Let $\mathcal{A}$ be a Grothendieck abelian category.
Let $\kappa$ be a cardinal as in
Lemma \ref{lemma-acyclic-quotient-complexes-bounded-size}.
Suppose that $I^\bullet$ is a complex such that
\begin{enumerate}
\item each $I^j$ is injective, and
\item for every bounded above acyclic complex $M^\bullet$
such that $|M^n| \leq \kappa$
we have $\Hom_{K(\mathcal{A})}(M^\bullet, I^\bullet) = 0$.
\end{enumerate}
Then $I^\bullet$ is an $K$-injective complex.
\end{lemma}

\begin{proof}
Let $M^\bullet$ be an acyclic complex. We are going to construct by
induction on the ordinal $\alpha$ an acyclic subcomplex
$K_\alpha^\bullet \subset M^\bullet$ as follows.
For $\alpha = 0$ we set $K_0^\bullet = 0$. For $\alpha > 0$
we proceed as follows:
\begin{enumerate}
\item If $\alpha = \beta + 1$ and $K_\beta^\bullet = M^\bullet$
then we choose $K_\alpha^\bullet = K_\beta^\bullet$.
\item If $\alpha = \beta + 1$ and $K_\beta^\bullet \not = M^\bullet$
then $M^\bullet/K_\beta^\bullet$ is a nonzero acyclic complex.
We choose a subcomplex $N_\alpha^\bullet \subset M^\bullet/K_\beta^\bullet$
as in Lemma \ref{lemma-acyclic-quotient-complexes-bounded-size}.
Finally, we let $K_\alpha^\bullet \subset M^\bullet$
be the inverse image of $N_\alpha^\bullet$.
\item If $\alpha$ is a limit ordinal we set
$K_\beta^\bullet = \colim K_\alpha^\bullet$.
\end{enumerate}
It is clear that $M^\bullet = K_\alpha^\bullet$ for a suitably large
ordinal $\alpha$. We will prove that
$$
\Hom_{K(\mathcal{A})}(K_\alpha^\bullet, I^\bullet)
$$
is zero by transfinite induction on $\alpha$. It holds for $\alpha = 0$
since $K_0^\bullet$ is zero. Suppose it holds for $\beta$ and
$\alpha = \beta + 1$. In case (1) of the list above the result is clear.
In case (2) there is a short exact sequence of complexes
$$
0 \to K_\beta^\bullet \to K_\alpha^\bullet \to N_\alpha^\bullet \to 0
$$
Since each component of $I^\bullet$ is injective we see that we obtain
an exact sequence
$$
\Hom_{K(\mathcal{A})}(K_\beta^\bullet, I^\bullet) \leftarrow
\Hom_{K(\mathcal{A})}(K_\alpha^\bullet, I^\bullet) \leftarrow
\Hom_{K(\mathcal{A})}(N_\alpha^\bullet, I^\bullet)
$$
By induction the term on the left is zero and by assumption on $I^\bullet$
the term on the right is zero. Thus the middle group is zero too.
Finally, suppose that $\alpha$ is a limit ordinal. Then we see that
$$
\Hom^\bullet(K_\alpha^\bullet, I^\bullet) =
\lim_{\beta < \alpha} \Hom^\bullet(K_\beta^\bullet, I^\bullet)
$$
with notation as in
More on Algebra, Section \href{https://stacks.math.columbia.edu/tag/0A8H}{section hom complexes (Tag 0A8H)}.
These complexes compute morphisms in $K(\mathcal{A})$ by
More on Algebra, Equation
(\href{https://stacks.math.columbia.edu/tag/0A5X}{equation cohomology hom complex (Tag 0A5X)}).
Note that the transition maps in the system are surjective
because $I^j$ is injective for each $j$. Moreover, for a limit
ordinal $\alpha$ we have equality of limit and value
(see displayed formula above). Thus we may apply
Homology, Lemma \href{https://stacks.math.columbia.edu/tag/0AAT}{lemma ML over ordinals (Tag 0AAT)}
to conclude.
\end{proof}


\begin{lemma}
\label{lemma-functorial-homotopies}
Let $\mathcal{A}$ be a Grothendieck abelian category.
Let $(K_i^\bullet)_{i \in I}$ be a set of acyclic complexes.
There exists a functor $M^\bullet \mapsto \mathbf{M}^\bullet(M^\bullet)$
and a natural transformation
$j_{M^\bullet} : M^\bullet \to \mathbf{M}^\bullet(M^\bullet)$
such
\begin{enumerate}
\item $j_{M^\bullet}$ is a (termwise) injective quasi-isomorphism, and
\item for every $i \in I$ and $w : K_i^\bullet \to M^\bullet$
the morphism $j_{M^\bullet} \circ w$ is homotopic to zero.
\end{enumerate}
\end{lemma}

\begin{proof}
For every $i \in I$ choose a (termwise) injective map of complexes
$K_i^\bullet \to L_i^\bullet$ which is homotopic to zero with
$L_i^\bullet$ quasi-isomorphic to zero. For example, take $L_i^\bullet$
to be the cone on the identity of $K_i^\bullet$.
We define $\mathbf{M}^\bullet(M^\bullet)$ by the following pushout diagram
$$
\xymatrix{
\bigoplus_{i \in I}
\bigoplus_{w : K_i^\bullet \to M^\bullet}
K_i^\bullet \ar[r] \ar[d] & M^\bullet \ar[d] \\
\bigoplus_{i \in I}
\bigoplus_{w : K_i^\bullet \to M^\bullet}
L_i^\bullet \ar[r] &  \mathbf{M}^\bullet(M^\bullet).
}
$$
Then $M^\bullet \mapsto \mathbf{M}^\bullet(M^\bullet)$ is a functor. The right
vertical arrow defines the functorial injective map $j_{M^\bullet}$.
The cokernel of $j_{M^\bullet}$ is isomorphic to the direct sum of
the cokernels of the maps $K_i^\bullet \to L_i^\bullet$ hence acyclic.
Thus $j_{M^\bullet}$ is a quasi-isomorphism. Part (2) holds by construction.
\end{proof}


\begin{lemma}
\label{lemma-functorial-injective}
Let $\mathcal{A}$ be a Grothendieck abelian category.
There exists a functor $M^\bullet \mapsto \mathbf{N}^\bullet(M^\bullet)$
and a natural transformation
$j_{M^\bullet} : M^\bullet \to \mathbf{N}^\bullet(M^\bullet)$
such
\begin{enumerate}
\item $j_{M^\bullet}$ is a (termwise) injective quasi-isomorphism, and
\item for every $n \in \mathbf{Z}$ the map $M^n \to \mathbf{N}^n(M^\bullet)$
factors through a subobject $I^n \subset \mathbf{N}^n(M^\bullet)$ where $I^n$
is an injective object of $\mathcal{A}$.
\end{enumerate}
\end{lemma}

\begin{proof}
Choose a functorial injective embeddings $i_M : M \to I(M)$, see
Theorem \ref{theorem-injective-embedding-grothendieck}.
For every complex $M^\bullet$ denote $J^\bullet(M^\bullet)$ the complex
with terms $J^n(M^\bullet) = I(M^n) \oplus I(M^{n + 1})$ and differential
$$
d_{J^\bullet(M^\bullet)} =
\left(
\begin{matrix}
0 & 1 \\
0 & 0
\end{matrix}
\right)
$$
There exists a canonical injective map of complexes
$u_{M^\bullet} : M^\bullet \to J^\bullet(M^\bullet)$ by mapping $M^n$ to
$I(M^n) \oplus I(M^{n + 1})$ via the maps $i_{M^n} : M^n \to I(M^n)$ and
$i_{M^{n + 1}} \circ d : M^n \to M^{n + 1} \to I(M^{n + 1})$. Hence a
short exact sequence of complexes
$$
0 \to M^\bullet \xrightarrow{u_{M^\bullet}}
J^\bullet(M^\bullet) \xrightarrow{v_{M^\bullet}}
Q^\bullet(M^\bullet) \to 0
$$
functorial in $M^\bullet$. Set
$$
\mathbf{N}^\bullet(M^\bullet) = C(v_{M^\bullet})^\bullet[-1].
$$
Note that
$$
\mathbf{N}^n(M^\bullet) = Q^{n - 1}(M^\bullet) \oplus J^n(M^\bullet)
$$
with differential
$$
\left(
\begin{matrix}
- d^{n - 1}_{Q^\bullet(M^\bullet)} & - v^n_{M^\bullet} \\
0 & d^n_{J^\bullet(M)}
\end{matrix}
\right)
$$
Hence we see that there is a map of complexes
$j_{M^\bullet} : M^\bullet \to \mathbf{N}^\bullet(M^\bullet)$
induced by $u$. It is injective and factors through an injective subobject
by construction. The map $j_{M^\bullet}$ is a quasi-isomorphism as one
can prove by looking at the long exact sequence of cohomology associated
to the short exact sequences of complexes above.
\end{proof}


\begin{theorem}
\label{theorem-injective-embedding-grothendieck}
Let $\mathcal{A}$ be a Grothendieck abelian category.
Then $\mathcal{A}$ has functorial injective embeddings.
\end{theorem}

\begin{proof}
Please compare with the proof of
Theorem \href{https://stacks.math.columbia.edu/tag/05NX}{theorem baer grothendieck (Tag 05NX)}.
Choose a generator $U$ of $\mathcal{A}$. For an object $M$ we define
$\mathbf{M}(M)$ by the following pushout diagram
$$
\xymatrix{
\bigoplus_{N \subset U}
\bigoplus_{\varphi \in \Hom(N, M)}
N \ar[r] \ar[d] & M \ar[d] \\
\bigoplus_{N \subset U}
\bigoplus_{\varphi \in \Hom(N, M)}
U \ar[r] &  \mathbf{M}(M).
}
$$
Note that $M \mapsto \mathbf{M}(M)$ is a functor and that there
exist functorial injective maps $M \to \mathbf{M}(M)$. By transfinite
induction we define functors $\mathbf{M}_\alpha(M)$ for every
ordinal $\alpha$. Namely, set $\mathbf{M}_0(M) = M$. Given
$\mathbf{M}_\alpha(M)$ set
$\mathbf{M}_{\alpha + 1}(M) = \mathbf{M}(\mathbf{M}_\alpha(M))$.
For a limit ordinal $\beta$ set
$$
\mathbf{M}_\beta(M) = \colim_{\alpha < \beta} \mathbf{M}_\alpha(M).
$$
Finally, pick any ordinal $\alpha$ whose cofinality is greater than $|U|$.
Such an ordinal exists by
Sets, Proposition \href{https://stacks.math.columbia.edu/tag/05N3}{proposition exist ordinals large cofinality (Tag 05N3)}.
We claim that $M \to \mathbf{M}_\alpha(M)$ is the desired functorial
injective embedding. Namely, if $N \subset U$ is a subobject and
$\varphi : N \to \mathbf{M}_\alpha(M)$ is a morphism, then we see that
$\varphi$ factors through $\mathbf{M}_{\alpha'}(M)$ for some
$\alpha' < \alpha$ by
Proposition \ref{proposition-objects-are-small}.
By construction of $\mathbf{M}(-)$ we see that $\varphi$ extends to
a morphism from $U$ into $\mathbf{M}_{\alpha' + 1}(M)$ and hence into
$\mathbf{M}_\alpha(M)$. By
Lemma \ref{lemma-characterize-injective}
we conclude that $\mathbf{M}_\alpha(M)$ is injective.
\end{proof}


\begin{lemma}
\label{lemma-characterize-injective}
\begin{slogan}
To check that an object is injective, one only needs to check that lifting
holds for subobjects of a generator.
\end{slogan}
Let $\mathcal{A}$ be a Grothendieck abelian category with generator $U$.
An object $I$ of $\mathcal{A}$ is injective if and only if in every
commutative diagram
$$
\xymatrix{
M \ar[d] \ar[r] &  I \\
U \ar@{-->}[ru]
}
$$
for $M \subset U$ a subobject, the dotted arrow exists.
\end{lemma}

\begin{proof}
Please see Lemma \href{https://stacks.math.columbia.edu/tag/05NU}{lemma criterion baer (Tag 05NU)} for the case of modules.
Choose an injection $A \subset B$ and a morphism $\varphi : A \to I$.
Consider the set $S$ of pairs $(A', \varphi')$ consisting of
subobjects $A \subset A' \subset B$ and a morphism $\varphi' : A' \to I$
extending $\varphi$. Define a partial ordering on this set in the obvious
manner. Choose a totally ordered subset $T \subset S$. Then
$$
A' = \colim_{t \in T} A_t \xrightarrow{\colim_{t \in T} \varphi_t} I
$$
is an upper bound. Hence by Zorn's lemma the set $S$ has a maximal element
$(A', \varphi')$. We claim that $A' = B$. If not, then choose a morphism
$\psi : U \to B$ which does not factor through $A'$. Set
$N = A' \cap \psi(U)$. Set $M = \psi^{-1}(N)$. Then the map
$$
M \to N \to A' \xrightarrow{\varphi'} I
$$
can be extended to a morphism $\chi : U \to I$. Since
$\chi|_{\Ker(\psi)} = 0$ we see that $\chi$ factors as
$$
U \to \Im(\psi) \xrightarrow{\varphi''} I
$$
Since $\varphi'$ and $\varphi''$ agree on $N = A' \cap \Im(\psi)$
we see that combined the define a morphism $A' + \Im(\psi) \to I$
contradicting the assumed maximality of $A'$.
\end{proof}


\begin{proposition}
\label{proposition-objects-are-small}
Let $\mathcal{A}$ be a Grothendieck abelian category. Let $M$ be an
object of $\mathcal{A}$. Let $\kappa = |M|$.
If $\alpha$ is an ordinal whose cofinality is bigger than $\kappa$,
then $M$ is $\alpha$-small with respect to injections.
\end{proposition}

\begin{proof}
Please compare with Proposition \href{https://stacks.math.columbia.edu/tag/05NT}{proposition modules are small (Tag 05NT)}.
We need only show that the map (\ref{equation-compare}) is a surjection.
Let $f : M \to \colim B_\beta$ be a map.
Consider the subobjects $\{f^{-1}(B_\beta)\}$ of $M$, where $B_\beta$
is considered as a subobject of the colimit $B = \bigcup_\beta B_\beta$.
If one of these, say $f^{-1}(B_\beta)$, fills $M$,
then the map factors through $B_\beta$.

\medskip\noindent
Because $\mathcal{A}$ has AB5 we have
$$
\colim f^{-1}(B_\beta) = f^{-1}\left(\colim B_\beta\right) = M.
$$
Now there are at most $\kappa$ different subobjects of $M$ that occur among
the $f^{-1}(B_\alpha)$, by hypothesis.
Thus we can find a subset $S \subset \alpha$ of cardinality at most
$\kappa$ such that as $\beta'$ ranges over $S$, the
$f^{-1}(B_{\beta'})$ range over \emph{all} the $f^{-1}(B_\alpha)$.

\medskip\noindent
However, $S$ has an upper bound $\widetilde{\alpha} < \alpha$ as
$\alpha$ has cofinality bigger than $\kappa$. In particular, all the
$f^{-1}(B_{\beta'})$, $\beta' \in S$ are contained in
$f^{-1}(B_{\widetilde{\alpha}})$.
It follows that $f^{-1}(B_{\widetilde{\alpha}}) = M$.
In particular, the map $f$ factors through $B_{\widetilde{\alpha}}$.
\end{proof}


\begin{lemma}
\label{lemma-set-iso-classes-bounded-size}
Let $\mathcal{A}$ be a Grothendieck abelian category with generator $U$.
\begin{enumerate}
\item If $|M| \leq \kappa$, then $M$ is the quotient of a direct
sum of at most $\kappa$ copies of $U$.
\item For every cardinal $\kappa$ the isomorphism classes
of objects $M$ with $|M| \leq \kappa$ form a set.
\end{enumerate}
\end{lemma}

\begin{proof}
For (1) choose for every proper subobject $M' \subset M$ a morphism
$\varphi_{M'} : U \to M$ whose image is not contained in $M'$. Then
$\bigoplus_{M' \subset M} \varphi_{M'} : \bigoplus_{M' \subset M} U \to M$
is surjective. It is clear that (1) implies (2).
\end{proof}


\begin{lemma}
\label{lemma-set-of-subobjects}
Let $\mathcal{A}$ be an abelian category with a generator $U$ and
$X$ an object of $\mathcal{A}$. If $\kappa$ is the cardinality of
$\Mor(U, X)$ then
\begin{enumerate}
\item There does not exist a strictly increasing
(or strictly decreasing) chain of subobjects
of $X$ indexed by a cardinal bigger than $\kappa$.
\item If $\alpha$ is an ordinal of cofinality $> \kappa$
then any increasing (or decreasing) sequence of subobjects
of $X$ indexed by $\alpha$ is eventually constant.
\item The cardinality of the set of subobjects of $X$
is $\leq 2^\kappa$.
\end{enumerate}
\end{lemma}

\begin{proof}
For (1) assume $\kappa' > \kappa$ is a cardinal and assume
$X_i$, $i \in \kappa'$ is strictly increasing. Then take for
each $i$ a $\phi_i \in \Mor(U, X)$ such that $\phi_i$ factors through
$X_{i + 1}$ but not through $X_i$. Then the morphisms $\phi_i$
are distinct, which contradicts the definition of $\kappa$.

\medskip\noindent
Part (2) follows from the definition of cofinality and (1).

\medskip\noindent
Proof of (3). For any subobject $Y \subset X$
define $S_Y \in \mathcal{P}(\Mor(U, X))$ (power set) as
$S_Y = \{\phi \in \Mor(U,X) : \phi)\text{ factors through }Y\}$.
Then $Y = Y'$ if and only if $S_Y = S_{Y'}$. Hence the cardinality
of the set of subobjects is at most the cardinality of this power set.
\end{proof}


\begin{theorem}
\label{theorem-K-injective-embedding-grothendieck}
\begin{slogan}
Existence of K-injective complexes for Grothendieck abelian categories.
\end{slogan}
Let $\mathcal{A}$ be a Grothendieck abelian category.
For every complex $M^\bullet$ there exists a quasi-isomorphism
$M^\bullet \to I^\bullet$ such that $M^n \to I^n$ is injective and $I^n$
is an injective object of $\mathcal{A}$ for all $n$ and $I^\bullet$
is a K-injective complex. Moreover, the construction is functorial in
$M^\bullet$.
\end{theorem}

\begin{proof}
Please compare with the proof of
Theorem \href{https://stacks.math.columbia.edu/tag/05NX}{theorem baer grothendieck (Tag 05NX)}
and
Theorem \ref{theorem-injective-embedding-grothendieck}.
Choose a cardinal $\kappa$ as in
Lemmas \ref{lemma-acyclic-quotient-complexes-bounded-size} and
\ref{lemma-characterize-K-injective}.
Choose a set $(K_i^\bullet)_{i \in I}$
of bounded above, acyclic complexes
such that every bounded above acyclic complex $K^\bullet$
such that $|K^n| \leq \kappa$ is isomorphic to $K_i^\bullet$ for some
$i \in I$. This is possible by
Lemma \ref{lemma-set-iso-classes-bounded-size}.
Denote $\mathbf{M}^\bullet(-)$ the functor constructed in
Lemma \ref{lemma-functorial-homotopies}.
Denote $\mathbf{N}^\bullet(-)$ the functor constructed in
Lemma \ref{lemma-functorial-injective}.
Both of these functors come with injective transformations
$\text{id} \to \mathbf{M}$ and $\text{id} \to \mathbf{N}$.

\medskip\noindent
Using transfinite recursion we define a sequence of functors
$\mathbf{T}_\alpha(-)$ and corresponding transformations
$\text{id} \to \mathbf{T}_\alpha$. Namely we set
$\mathbf{T}_0(M^\bullet) = M^\bullet$. If $\mathbf{T}_\alpha$ is
given then we set
$$
\mathbf{T}_{\alpha + 1}(M^\bullet) =
\mathbf{N}^\bullet(\mathbf{M}^\bullet(\mathbf{T}_\alpha(M^\bullet)))
$$
If $\beta$ is a limit ordinal we set
$$
\mathbf{T}_\beta(M^\bullet) =
\colim_{\alpha < \beta} \mathbf{T}_\alpha(M^\bullet)
$$
The transition maps of the system are injective quasi-isomorphisms.
By AB5 we see that the colimit is still quasi-isomorphic to $M^\bullet$.
We claim that $M^\bullet \to \mathbf{T}_\alpha(M^\bullet)$
does the job if the cofinality of $\alpha$ is larger than
$\max(\kappa, |U|)$ where $U$ is a generator of $\mathcal{A}$.
Namely, it suffices to check conditions (1) and (2) of
Lemma \ref{lemma-characterize-K-injective}.

\medskip\noindent
For (1) we use the criterion of
Lemma \ref{lemma-characterize-injective}.
Suppose that $M \subset U$ and $\varphi : M \to \mathbf{T}^n_\alpha(M^\bullet)$
is a morphism for some $n \in \mathbf{Z}$. By
Proposition \ref{proposition-objects-are-small}
we see that $\varphi$ factor through
$\mathbf{T}^n_{\alpha'}(M^\bullet)$ for some $\alpha' < \alpha$.
In particular, by the construction of the functor
$\mathbf{N}^\bullet(-)$ we see that $\varphi$ factors through
an injective object of $\mathcal{A}$ which shows that $\varphi$
lifts to a morphism on $U$.

\medskip\noindent
For (2) let $w : K^\bullet  \to \mathbf{T}_\alpha(M^\bullet)$
be a morphism of complexes where $K^\bullet$ is a bounded above acyclic
complex such that $|K^n| \leq \kappa$. Then $K^\bullet \cong K_i^\bullet$
for some $i \in I$. Moreover, by
Proposition \ref{proposition-objects-are-small}
once again we see that $w$ factor through
$\mathbf{T}^n_{\alpha'}(M^\bullet)$ for some $\alpha' < \alpha$.
In particular, by the construction of the functor
$\mathbf{M}^\bullet(-)$ we see that $w$ is homotopic to zero.
This finishes the proof.
\end{proof}
\end{document}
